% !TEX root = ../main.tex
% Task 2B: Winning classification algorithms in a real competition.
% Required items: competition description, winner + technique, main idea, what makes it stand out.

\section{Winning classification algorithms}\label{sec:winning-classification}

The Otto Group Product Classification Challenge ran on Kaggle from March to May 2015 \cite{kaggle2015otto}. Competitors had to assign each of 144\,368 unlabelled products to one of nine commercial categories using 61\,878 labelled training samples described by 93 anonymised integer features. Submissions were scored by multi-class logarithmic loss, which rewards calibrated probability estimates rather than hard labels. The contest attracted 3\,514 teams.

\paragraph{Winner and technique.} First place went to Gilberto Titericz and Stanislav Semenov (handle \textit{tin rtyz}). Their submission was a weighted three-layer stacked ensemble of about thirty base learners \cite{otto2015writeup}: gradient-boosted trees (XGBoost with \texttt{multi:softprob} under varied feature subsets and depths), bagged feedforward neural networks built in Keras, random forests and extra-trees classifiers, and a k-nearest-neighbour model whose distances were also added as extra features. Stratified out-of-fold predictions from these models fed a second stacking layer, and a conservative third-layer blender (logistic regression or a shallow XGBoost) produced the final probabilities.

\paragraph{Main idea.} The winners argue that diversity matters more than the quality of any single model. They deliberately included weak but uncorrelated learners, such as KNN which was not competitive on its own, because their errors cancel inside the stack. The theoretical basis is Wolpert's stacked generalisation \cite{wolpert1992stacked}: blending out-of-fold predictions from heterogeneous models reduces variance without adding bias, as long as the meta-learner never sees in-fold predictions. XGBoost \cite{chen2016xgboost} supplied the gradient-boosted component and absorbed the 93 anonymous features efficiently thanks to its regularised objective and histogram-based split finding.

\paragraph{What makes it stand out.} Three choices distinguish the winning entry from a typical competitive submission. First, every layer was trained with stratified 5-fold cross-validation and the out-of-fold predictions were tracked carefully so the meta-learner never saw leakage, a discipline most single-model entries skipped. Second, the neural networks were bagged across bootstrap samples, which stabilised an architecture that is otherwise variance-prone. Third, KNN was used as a source of meta-features rather than as a standalone predictor, capturing local geometric structure that tree ensembles smooth over. The combined effect was a public-leaderboard log-loss of about $0.38$, against roughly $0.44$ for a well-tuned single XGBoost (top 10\%) and $0.46$ for a single Keras network (top 15\%). A plain average of those two single models already reaches the top 3\%; the full three-layer stack adds another $0.03$, which was enough to win against several thousand competitors.
